% Figura: ley de Coulomb (geometría) y principio de superposición.
\begin{figure}[ht]
\centering
\tdplotsetmaincoords{72}{125}
\begin{minipage}[t]{0.49\textwidth}\centering
\begin{tikzpicture}[tdplot_main_coords, scale=1.0, baseline=(current bounding box.north)]
  \draw[eje] (0,0,0) -- (3.2,0,0) node[below left] {\small $x$};
  \draw[eje] (0,0,0) -- (0,3.4,0) node[right] {\small $y$};
  \draw[eje] (0,0,0) -- (0,0,3.2) node[above] {\small $z$};
  \coordinate (Q1) at (2.6,0.4,0.6);
  \coordinate (Q2) at (0.4,2.7,2.3);
  \draw[->, azulclaro!80!azulacero, line width=0.7pt] (0,0,0) -- (Q1) node[pos=0.5, below, font=\small, text=azulprof] {$\vc{r}_1$};
  \draw[->, azulclaro!80!azulacero, line width=0.7pt] (0,0,0) -- (Q2) node[pos=0.45, left, font=\small, text=azulprof] {$\vc{r}_2$};
  \draw[vect, azulprof, line width=1.2pt] (Q1) -- (Q2) node[pos=0.62, above left, font=\small, text=azulprof] {$\vc{R}_{12}$};
  \draw[vectrojo] (Q2) -- ($(Q2)+(-0.45,0.7,0.65)$) node[above, font=\small, text=acento] {$\vc{F}_{12}$};
  \draw[vectrojo] (Q1) -- ($(Q1)+(0.45,-0.7,-0.65)$) node[left, font=\small, text=acento] {$\vc{F}_{21}$};
  \node[cargapos] at (Q1) {$+$};
  \node[cargapos] at (Q2) {$+$};
  \node[font=\small, anchor=south west] at ($(Q1)+(0.1,0.1,0.1)$) {$q_1$};
  \node[font=\small, anchor=south east] at ($(Q2)+(-0.1,0,0.15)$) {$q_2$};
\end{tikzpicture}\\[2pt]{\small\scshape (a)}
\end{minipage}\hfill
\begin{minipage}[t]{0.49\textwidth}\centering
\begin{tikzpicture}[tdplot_main_coords, scale=0.95, baseline=(current bounding box.north)]
  \draw[eje] (0,0,0) -- (3.0,0,0) node[below left] {\small $x$};
  \draw[eje] (0,0,0) -- (0,3.2,0) node[right] {\small $y$};
  \draw[eje] (0,0,0) -- (0,0,3.0) node[above] {\small $z$};
  \coordinate (P) at (1.3,1.6,1.1);
  \foreach \n/\a/\b/\c/\l in {1/2.9/0.0/0.3/{q_1}, 2/2.5/1.5/2.4/{q_2}, 3/0.2/3.1/1.8/{q_3}}{
     \coordinate (S\n) at (\a,\b,\c);
     \draw[vect, azulclaro!90!black, line width=0.7pt] (S\n) -- (P);
     \node[cargapos] at (S\n) {$+$};
     \node[font=\small, anchor=south east] at ($(S\n)+(-0.1,0,0.1)$) {$\l$};}
  \draw[vectrojo, line width=1.3pt] (P) -- ($(P)+(-0.9,0.6,0.45)$) node[above left, font=\small, text=acento] {$\vc{F}_{q}$};
  \node[cargapru] at (P) {$q$};
  \node[font=\scriptsize, text=azulprof, anchor=north] at ($(S1)!0.5!(P)$) {$\vc{R}_1$};
  \node[font=\scriptsize, text=azulprof, anchor=east] at ($(S2)!0.5!(P)$) {$\vc{R}_2$};
  \node[font=\scriptsize, text=azulprof, anchor=west] at ($(S3)!0.45!(P)$) {$\vc{R}_3$};
\end{tikzpicture}\\[2pt]{\small\scshape (b)}
\end{minipage}
\caption{(a)~Ley de Coulomb para dos cargas puntuales del mismo signo:
$\vc{R}_{12}=\vc{r}_2-\vc{r}_1$ apunta de $q_1$ a $q_2$, y la fuerza
$\vc{F}_{12}$ sobre $q_2$ es paralela a $\vc{R}_{12}$ (repulsión); la fuerza sobre $q_1$ es la
opuesta, $\vc{F}_{21}=-\vc{F}_{12}$. (b)~Superposición: la fuerza sobre la carga $q$ es la suma
vectorial de las ejercidas por cada $q_i$, con $\vc{R}_i$ el vector que une $q_i$ con $q$.}
\label{fig:coulomb}
\end{figure}
